% English draft based on sections/10_gettier_problem_ja.tex
\section{The Gettier Problem: Stability Is Not Truth}

The Gettier problem shows that justified true belief can still fail to be knowledge. In the present framework, this is not an anomaly in epistemology. It is evidence that closure degree $C$ is not a truth measure.

\subsection{Structural Diagnosis}

A belief can be justified, true, and yet accidentally true. Its stopping structure is stable enough to be accepted, but the path by which it arrived at truth is not structurally appropriate. The problem is therefore not simply missing justification. It is the non-identity of stability and truth.

In this paper, $L_E$ supplies revisable stopping, not guaranteed truth. A stopping point can be operationally stable while still requiring later correction. This is precisely why intelligence needs logs, external reference fields, and revision mechanisms.

\subsection{$C$ as Stability Measure}

Closure degree $C$ measures operable stability. It does not measure truth. If $C$ were truth, then higher closure would monotonically approach knowledge. Gettier cases show that this is false. A state can possess high closure and still fail as knowledge.

The correct reading is:
\[
C = \text{degree of operational closure/stability},
\]
not
\[
C = \text{truth}.
\]

\subsection{Error Tolerance and $\mu_\theta$}

The ability to move thresholds entails the ability to be wrong. If a system can reorganize its stopping region, it can install a stopping operator in a region that later proves inadequate. Error is not an accidental defect added to intelligence from outside. It is structurally linked to threshold mobility.

A perfectly correct intelligence would never need to move thresholds. But a system that never moves thresholds is rigid. Thus error tolerance and $\mu_\theta$ are structurally coupled.

\subsection{Consequence}

Knowledge requires more than immediate stopping. It requires revisable stopping supported by logs, external reference fields, and later correction. The Gettier problem therefore confirms the paper's central claim: intelligence persists not by reaching absolute closure, but by operating provisional stopping under conditions of possible revision.

\bigskip
